\section{\normalsize Propositional Logic}
\vspace{-0.5em}
{\color{blue}\subsection{Propositions}}
A proposition is a declarative statement that has exactly one truth value: either True ($T$) or False ($F$).
{\color{blue}\subsection{Logical Connectives}}
\subsubsection{Logical Operators}
$$
\begingroup
\scriptsize
\renewcommand{\arraystretch}{1.2}
\noindent\resizebox{\linewidth}{!}{$
\begin{array}{|c|c|c|}
\hline
\textbf{Symbol} & \textbf{Name} & \textbf{Meaning}\\
\hline
\neg p & \text{Negation} & \text{NOT }p\\
\hline
p\land q & \text{Conjunction} & p\text{ AND }q\\
\hline
p\lor q & \text{Disjunction} & p\text{ OR }q\\
\hline
p\to q & \text{Implication} & \text{IF }p\text{ THEN }q\\
\hline
p\leftrightarrow q & \text{Biconditional} & \text{IFF}\\
\hline
p\oplus q & \text{XOR} & \text{Exactly one}\\
\hline
\end{array}
$}
\endgroup
$$
\subsubsection{Truth Table}
$$
\begingroup
\scriptsize
\renewcommand{\arraystretch}{1.2}
\noindent\resizebox{\linewidth}{!}{$
\begin{array}{|c|c|c|c|c|c|c|c|}
\hline
p&q&\neg p&p\land q&p\lor q&p\to q&p\leftrightarrow q&p\oplus q\\
\hline
T&T&F&T&T&T&T&F\\
\hline
T&F&F&F&T&F&F&T\\
\hline
F&T&T&F&T&T&F&T\\
\hline
F&F&T&F&F&T&T&F\\
\hline
\end{array}
$}
\endgroup
$$
{\color{blue}\subsection{Propositional \\ Equivalences}}
\subsubsection{Logical Equivalences}
$$
\begingroup
\scriptsize
\renewcommand{\arraystretch}{1.2}
\noindent\resizebox{\linewidth}{!}{$
\begin{array}{|l|l|}
\hline
\text{Law} & \text{Equivalence}\\
\hline
\text{Identity} &
p\land T\equiv p,\quad p\lor F\equiv p\\
\hline
\text{Domination} &
p\lor T\equiv T,\quad p\land F\equiv F\\
\hline
\text{Idempotent} &
p\land p\equiv p,\quad p\lor p\equiv p\\
\hline
\text{Double Negation} &
\neg\neg p\equiv p\\
\hline
\text{Commutative} &
p\land q\equiv q\land p\\
& p\lor q\equiv q\lor p\\
\hline
\text{Associative} &
(p\land q)\land r\equiv p\land(q\land r)\\
& (p\lor q)\lor r\equiv p\lor(q\lor r)\\
\hline
\text{Distributive} &
p\land(q\lor r)\equiv(p\land q)\lor(p\land r)\\
& p\lor(q\land r)\equiv(p\lor q)\land(p\lor r)\\
\hline
\text{De Morgan} &
\neg(p\land q)\equiv\neg p\lor\neg q\\
& \neg(p\lor q)\equiv\neg p\land\neg q\\
\hline
\text{Negation} &
p\lor\neg p\equiv T,\quad p\land\neg p\equiv F\\
\hline
\text{Absorption} &
p\lor(p\land q)\equiv p\\
& p\land(p\lor q)\equiv p\\
\hline
\end{array}
$}
\endgroup
$$
$$
\begingroup
\scriptsize
\renewcommand{\arraystretch}{1.2}
\noindent\resizebox{\linewidth}{!}{$
\begin{array}{|l|l|}
\hline
\text{Law} & \text{Equivalence}\\
\hline
\text{Conditional} & p\to q\equiv\neg p\lor q\\
\hline
\text{Converse} & q\to p\\
\hline
\text{Inverse} & \neg p\to\neg q\\
\hline
\text{Contrapositive} & \neg q\to\neg p\\
\hline
\text{Negation of Conditional} &
\neg(p\to q)\equiv p\land\neg q\\
\hline
\end{array}
$}
\endgroup
$$
\subsubsection{Logical Equivalences \\ Involving Conditional and Biconditional \\ Statements}
$$
\begingroup
\scriptsize
\renewcommand{\arraystretch}{1.2}
\noindent\resizebox{\linewidth}{!}{$
\begin{array}{|l|l|}
\hline
\text{Law} & \text{Equivalence}\\
\hline
\text{Conditional Identity} &
p\to q\equiv\neg p\lor q\\
\hline
\text{Contrapositive} &
p\to q\equiv\neg q\to\neg p\\
\hline
\text{Negation of Conditional} &
\neg(p\to q)\equiv p\land\neg q\\
\hline
\text{Biconditional Identity} &
p\leftrightarrow q\equiv(p\to q)\land(q\to p)\\
\hline
\text{Biconditional (Expanded)} &
p\leftrightarrow q\equiv
(p\land q)\lor(\neg p\land\neg q)\\
\hline
\text{Negation of Biconditional} &
\neg(p\leftrightarrow q)\equiv
(p\land\neg q)\lor(\neg p\land q)\\
\hline
\end{array}
$}
\endgroup
$$
\subsubsection{Translation of Logical \\ Equivalences}
$$
\begingroup
\scriptsize
\renewcommand{\arraystretch}{1.2}
\noindent\resizebox{\linewidth}{!}{$
\begin{array}{|l|l|}
\hline
\text{English Phrase} & \text{Logical Form}\\
\hline
\text{not }p & \neg p\\
\hline
p\text{ and }q & p\land q\\
\hline
p\text{ or }q & p\lor q\\
\hline
\text{Either }p\text{ or }q\text{, but not both} & p\oplus q\\
\hline
\text{If }p\text{, then }q & p\to q\\
\hline
p\text{ only if }q & p\to q\\
\hline
p\text{ if }q & q\to p\\
\hline
p\text{ if and only if }q & p\leftrightarrow q\\
\hline
p\text{ is sufficient for }q & p\to q\\
\hline
p\text{ is necessary for }q & q\to p\\
\hline
\text{Neither }p\text{ nor }q & \neg p\land\neg q\\
\hline
\text{Not both }p\text{ and }q & \neg(p\land q)\\
\hline
\text{At least one of }p,q & p\lor q\\
\hline
\text{Both }p\text{ and }q\text{ are false} & \neg p\land\neg q\\
\hline
\end{array}
$}
\endgroup
$$
{\color{blue}\subsection{Predicates and \\ Quantifiers}}
A predicate $P(x)$ is a statement involving a variable
that becomes a proposition when a value is assigned to the variable.

The domain of discourse is the set of values a variable
is allowed to take.
$$
\begingroup
\scriptsize
\renewcommand{\arraystretch}{1.2}
\noindent\resizebox{\columnwidth}{!}{$
\begin{array}{|c|l|l|}
\hline
\text{Symbol} & \text{Name} & \text{Meaning}\\
\hline
\forall x\,P(x) & \text{Universal} &
P(x)\text{ holds for every }x\\
\hline
\exists x\,P(x) & \text{Existential} &
P(x)\text{ holds for at least one }x\\
\hline
\exists!x\,P(x) & \text{Unique existence} &
P(x)\text{ holds for exactly one }x\\
\hline
\end{array}
$}
\endgroup
$$
Example: Let $P(x)$ mean "$x>0$" over $\mathbb{Z}$.
Then $\forall x\,P(x)$ is false, while $\exists x\,P(x)$ is true.
\subsubsection{De Morgan's Laws for \\ Quantifiers}
$$
\begingroup
\scriptsize
\renewcommand{\arraystretch}{1.2}
\noindent\resizebox{\columnwidth}{!}{$
\begin{array}{|l|l|}
\hline
\text{Law} & \text{Equivalence}\\
\hline
\text{Universal Negation} &
\neg\forall x\,P(x)\equiv\exists x\,\neg P(x)\\
\hline
\text{Existential Negation} &
\neg\exists x\,P(x)\equiv\forall x\,\neg P(x)\\
\hline
\text{Negated Universal} &
\neg(\forall x\,P(x))\equiv\exists x\,\neg P(x)\\
\hline
\text{Negated Existential} &
\neg(\exists x\,P(x))\equiv\forall x\,\neg P(x)\\
\hline
\end{array}
$}
\endgroup
$$
\subsubsection{Multiple Quantifiers}
$$
\begingroup
\scriptsize
\renewcommand{\arraystretch}{1.2}
\noindent\resizebox{\columnwidth}{!}{$
\begin{array}{|l|l|}
\hline
\text{Statement} & \text{Meaning}\\
\hline
\forall x\forall y\,P(x,y) &
P\text{ holds for every pair }(x,y)\\
\hline
\exists x\exists y\,P(x,y) &
P\text{ holds for at least one pair}\\
\hline
\forall x\exists y\,P(x,y) &
\text{For every }x,\text{ some }y\text{ works}\\
\hline
\exists y\forall x\,P(x,y) &
\text{One }y\text{ works for every }x\\
\hline
\end{array}
$}
\endgroup
$$
\subsubsection{Translation into \\ Predicate Logic}
$$
\begingroup
\scriptsize
\renewcommand{\arraystretch}{1.2}
\noindent\resizebox{\columnwidth}{!}{$
\begin{array}{|l|l|}
\hline
\text{English} & \text{Logical Form}\\
\hline
\text{Every }x\text{ has property }P &
\forall x\,P(x)\\
\hline
\text{Some }x\text{ has property }P &
\exists x\,P(x)\\
\hline
\text{No }x\text{ has property }P &
\forall x\,\neg P(x)\\
\hline
\text{Not every }x\text{ has property }P &
\exists x\,\neg P(x)\\
\hline
\text{Everyone is }P\text{ and }Q &
\forall x\,(P(x)\land Q(x))\\
\hline
\text{Everyone is }P\text{ or }Q &
\forall x\,(P(x)\lor Q(x))\\
\hline
\text{Someone is }P\text{ and }Q &
\exists x\,(P(x)\land Q(x))\\
\hline
\text{Someone is }P\text{ or }Q &
\exists x\,(P(x)\lor Q(x))\\
\hline
\text{Every }P\text{ is }Q &
\forall x\,(P(x)\to Q(x))\\
\hline
\text{Some }P\text{ is }Q &
\exists x\,(P(x)\land Q(x))\\
\hline
\text{No }P\text{ is }Q &
\forall x\,(P(x)\to\neg Q(x))\\
\hline
\text{Some }P\text{ is not }Q &
\exists x\,(P(x)\land\neg Q(x))\\
\hline
\end{array}
$}
\endgroup
$$
{\color{blue}\subsection{Rules of Inference}}
$$
\begingroup
\scriptsize
\renewcommand{\arraystretch}{1.15}
\noindent\resizebox{\columnwidth}{!}{$
\begin{array}{|l|c|}
\hline
\text{Rule} & \text{Form}\\
\hline
\text{Modus Ponens (MP)} &
\begin{array}{l}
p\\
p\to q\\
\hline
\therefore q
\end{array}\\
\hline
\text{Modus Tollens (MT)} &
\begin{array}{l}
p\to q\\
\neg q\\
\hline
\therefore\neg p
\end{array}\\
\hline
\text{Hypothetical Syllogism (HS)} &
\begin{array}{l}
p\to q\\
q\to r\\
\hline
\therefore p\to r
\end{array}\\
\hline
\text{Disjunctive Syllogism (DS)} &
\begin{array}{l}
p\lor q\\
\neg p\\
\hline
\therefore q
\end{array}\\
\hline
\text{Addition} &
\begin{array}{l}
p\\
\hline
\therefore p\lor q
\end{array}\\
\hline
\text{Simplification} &
\begin{array}{l}
p\land q\\
\hline
\therefore p
\end{array}\\
\hline
\text{Conjunction} &
\begin{array}{l}
p\\
q\\
\hline
\therefore p\land q
\end{array}\\
\hline
\text{Resolution} &
\begin{array}{l}
p\lor q\\
\neg p\lor r\\
\hline
\therefore q\lor r
\end{array}\\
\hline
\text{Universal Instantiation (UI)} &
\begin{array}{l}
\forall x\,P(x)\\
\hline
\therefore P(c)
\end{array}\\
\hline
\text{Existential Instantiation (EI)} &
\begin{array}{l}
\exists x\,P(x)\\
\hline
\therefore P(c)
\end{array}\\
\hline
\text{Universal Generalization (UG)} &
\begin{array}{l}
P(c),\ c\text{ arbitrary}\\
\hline
\therefore\forall x\,P(x)
\end{array}\\
\hline
\text{Existential Generalization (EG)} &
\begin{array}{l}
P(c)\\
\hline
\therefore\exists x\,P(x)
\end{array}\\
\hline
\end{array}
$}
\endgroup
$$
{\color{blue}\subsection{Introduction to Proofs}}
\subsubsection{Choosing \\ a Proof Method}
$$
\begingroup
\renewcommand{\arraystretch}{1.2}
\resizebox{\columnwidth}{!}{$
\begin{array}{|l|}
\hline
\text{Can }p\text{ directly imply }q\text{?}\\
\text{Direct Proof}\\
\hline
\text{Is }\neg q\to\neg p\text{ easier to prove?}\\
\text{Contraposition}\\
\hline
\text{Is the hypothesis always false?}\\
\text{Vacuous Proof}\\
\hline
\text{Is the conclusion always true?}\\
\text{Trivial Proof}\\
\hline
\text{Does assuming the opposite lead to a contradiction?}\\
\text{Proof by Contradiction}\\
\hline
\text{Can the statement be divided into exhaustive cases?}\\
\text{Proof by Cases}\\
\hline
\text{Is the statement about all integers }n\geq n_0\text{?}\\
\text{Mathematical Induction}\\
\hline
\text{Can one example disprove the claim?}\\
\text{Counterexample}\\
\hline
\end{array}
$}
\endgroup
$$